% Fragmento público generado desde propietarios íntegros.
% Residencia: 52.2.3.
% Fuente: ../incoming/radion_monografia_20260827/manuscrito/sections/03_cierre_exacto.tex:104-253
\subsubsection{Cociclo de transición entre las secciones directa y torsional}

\paragraph{Definición operatoria del acarreo \(\varepsilon_R\)}

\begin{definition}[Acarreo real del radión]
La salida unitaria del radión es el cociclo de transición
\begin{equation}
\epsone:=S_T-S_E
=\Phi_T-D_E
=\frac{20}{81}\Delta_4-operatorname{rem}_1(S_E).
\label{eq:epsilon-def}
\end{equation}
\end{definition}

La ecuación contiene una resta, pero no es un ajuste al blanco. Sus dos
operandos \(S_T\) y \(S_E\) se generan separadamente. La independencia que
se demuestra es independencia respecto de \(180/\pi\) y de un
\(\varepsilon\) objetivo; no se afirma independencia algebraica entre una
diferencia y sus sumandos.

\begin{theorem}[Cierre exacto unitario]
Con \(S_E\), \(\Phi_T\) y \(\epsone\) definidos por
\eqref{eq:seccion-electrica}, \eqref{eq:delta4}, \eqref{eq:rho-twoway} y
\eqref{eq:epsilon-def}, se cumple exactamente
\begin{equation}
\boxed{1=S_E-\Phi_T+\epsone.}
\label{eq:cierre-unitario}
\end{equation}
\end{theorem}
\begin{proof}
Por \eqref{eq:SE-intervalo}, \(D_E=S_E-1\). Por definición,
\(\epsone=\Phi_T-D_E\). Entonces
\[
S_E-\Phi_T+\epsone
=S_E-\Phi_T+\Phi_T-(S_E-1)=1.
\]
La fuerza de la prueba no está en esta cancelación elemental, sino en que la
genealogía previa fija los dos operandos sin usar el cierre como entrada.
\end{proof}

\paragraph{Publicación del cierre en la carta angular}

La carta de la vuelta multiplica una fracción unitaria por
\(360/T_+=180/\pih\). En particular,
\[
\epsR^\circ=\frac{360^\circ}{T_+}\epsone.
\]
Multiplicar \eqref{eq:cierre-unitario} por \(360^\circ/T_+\) y usar
\[
\frac{360}{T_+}S_E
=360\left(\frac5{36}+\frac{25}{9}\alphah\right)
=50+1000\alphah=50+A
\]
produce el cierre angular.

\begin{theorem}[Ecuación exacta del radión]
\BigRadionEquation
\end{theorem}

La ecuación dice algo más preciso que \enquote{el radián es
\(57.2957\ldots{}^\circ\)}. Descompone esa publicación en cuatro operaciones
con demostración de referencia:
\begin{center}
\begin{tabularx}{.96\textwidth}{>{\bfseries}l X}
\toprule
\(50^\circ\) & fase segunda de la rueda; bisagra crítica en una carta declarada;\\
\(A^\circ\) & coordenada común de la publicación parangular;\\
\(-\frac{20}{81}\Delta_4^\circ\) & transporte orientado de la torsión global;\\
\(+\epsR^\circ\) & acarreo real que empalma los dos lifts.\\
\bottomrule
\end{tabularx}
\end{center}

\paragraph{Convergencia por profundidad del refinamiento}

A profundidad \(N\), el generador entrega cilindros racionales
\(p_N\to\pih\) y \(e_N\to e_{\HMT}\). Se definen
\[
S_{E,N}=2p_N\left(\frac5{36}+\frac{25}{9}\alphah\right),
\]
\[
\Delta_{4,N}=\frac{p_N-e_N\log p_N}{270}
\left(1+\frac{p_N}{729}\right),\qquad
S_{T,N}=1+\frac{20}{81}\Delta_{4,N},
\]
y
\[
\mathcal E_N=S_{T,N}-S_{E,N}.
\]
La continuidad del funcional en \(3<p<4\), \(2<e<3\), junto a las
derivadas monótonas, transporta cada par de cilindros a un intervalo
certificado. A profundidad \(45\) bloques, su anchura es menor que
\(4.81\times10^{-130}\).

Cada retorno de nueve niveles contiene seis trits por nivel. Por eso la tasa
de refinamiento es
\begin{equation}
729^{-9}=3^{-54}.
\label{eq:contraccion-54}
\end{equation}
Esta cifra controla la anchura de los intervalos; no es el valor de
\(\epsR\).

La resta de tríadas estabiliza
\begin{center}
\ttfamily
000|000|000|896|802|822|044|562|981|999|050|695|020|651|286|413
\end{center}
y el prefijo de carries
\begin{center}
\ttfamily
0,0,0,0,0,-1,0,-1,-1,-1,0,-1,0.
\end{center}
Éste es el mecanismo concreto de memoria decimal: no una cola elegida, sino
la recurrencia \eqref{eq:subacarreo} aplicada a operandos previamente
construidos.

\paragraph{Evaluación de los valores canónicos}

La salida coinductiva es
\begin{align}
\epsone
&=8.9680282204456298199905069502065128641372736205009\ldots
\times10^{-10},\\
\epsR^\circ
&=5.1383016758575282071685164622645947489908574400054\ldots
\times10^{-8}\,{}^\circ.
\label{eq:epsilon-valores}
\end{align}
La publicación completa resulta
\[
\frac{180^\circ}{\pih}
=57.2957795130823208767981548141051703324054724665643\ldots{}^\circ.
\]

\begin{exactbox}
\begin{align*}
50^\circ+A^\circ
&=57.29735256928380099728510547238066299956\ldots{}^\circ,\\
\frac{20}{81}\Delta_4^\circ
&=0.00157310758449687906223272996065728980465\ldots{}^\circ,\\
50^\circ+A^\circ-\frac{20}{81}\Delta_4^\circ
&=57.29577946169930411822287274242000570976\ldots{}^\circ,\\
\epsR^\circ
&=0.00000005138301675857528207168516462264594749\ldots{}^\circ,\\
\text{suma final}
&=57.29577951308232087679815481410517033241\ldots{}^\circ.
\end{align*}
\end{exactbox}

